\section{Experimental protocol}
\label{sec:protocol}

\paragraph{Models, benchmarks and signals}
Backbones are LLaVA-1.5-7B \citep{llava} in 4-bit, Qwen2-VL-2B-Instruct \citep{qwenvl}, and LLaVA decoded with
visual contrastive decoding \citep{vcd}. Channel 2 is OWLv2-base \citep{owlv2} with $\theta=0.15$. Benchmarks are
the POPE adversarial split \citep{pope} and the AMBER discriminative benchmark \citep{amber}; MME, GQA and
HallusionBench appear only where the detector does not apply (\S\ref{sec:noevidence}). Per-item signals were
extracted once and are released, so every analysis below runs on a CPU. The extended POPE and AMBER cells (items $1500$--$2999$ of the
POPE adversarial split and items $500$--$2999$ of the shuffled AMBER list) were extracted with the same prompts and detector query as the
earlier cells; in the AMBER extension, photographs larger than $2000$ pixels on the long side were shrunk before detection because the
detector's preprocessing of the largest ones ($54$ megapixels) exhausted host memory.

\paragraph{Cohorts and the unit of splitting}
The extraction streamed the first $N$ rows of each public benchmark. We recovered the image identifier of every
item from the benchmark metadata and verified the reconstruction against the cached gold labels, object phrases
and categories (all agree). Table~\ref{tab:cohort} reports what enters each cell. The POPE items are
six questions per image (three present, three absent objects), so there are $75$, $100$, $250$ and $500$ distinct
images in the four POPE cell sizes; the smaller POPE cells are subsets of the largest one's images, and POPE-3000 is the whole
adversarial split. AMBER contributes $3.1$ items per image over $957$ images. Folds are unions of whole images, drawn at
random with a third of the images per fold. The intra-image correlation of the model's errors is slightly \emph{negative}
($-0.01$ to $-0.09$): in POPE the balanced yes/no design makes an error on one question of an image mildly predictive of a
correct answer on another.

\begin{table}[t]
\centering
\caption{Cohort flow. ``Detector value'' is the number of items for which the detector returned a score; the
remaining items stay in every analysis, keep the channel-1 decision, and cannot be repaired. In POPE the
exclusions are object phrases the extraction script failed to parse (POPE contains the typo ``imange''); in AMBER
they are the attribute and relation queries, which name no object. ICC is the intra-image correlation of the
model's errors.}
\label{tab:cohort}
\small
\resizebox{\textwidth}{!}{\begin{tabular}{lrrrrrlcc}
\toprule
Setting & items & images & items/img & ICC & detector value & excluded from the grounded cohort & images/fold & $\mu$ (\%)\\
\midrule
POPE-1500, LLaVA & 1500 & 250 & 6.0 & $-$0.04 & 1500 & 0 unparsed object phrases & 83/83/84 & 18.1\\
POPE-adv, LLaVA & 450 & 75 & 6.0 & $-$0.05 & 444 & 6 unparsed object phrases & 25/25/25 & 17.3\\
POPE-adv, Qwen2-VL & 600 & 100 & 6.0 & $-$0.04 & 591 & 9 unparsed object phrases & 33/33/34 & 12.7\\
POPE-adv, LLaVA+VCD & 600 & 100 & 6.0 & $-$0.04 & 591 & 9 unparsed object phrases & 33/33/34 & 19.7\\
AMBER(d), LLaVA & 500 & 392 & 1.3 & $-$0.09 & 228 & 272 attribute/relation queries & 130/130/132 & 21.4\\
\bottomrule
\end{tabular}
}
\end{table}

\paragraph{Primary analysis and missing evidence}
The primary cohort contains all items. Items with no detector value receive zeroed grounding features and a
missing-evidence flag in the score, and a margin $m=-1$ so they are never repaired. A secondary analysis restricted to
items with a detector value appears in Appendix~\ref{app:grounded}. Excluding the items would be an analysis of the
detector's coverage, not of deployment: a system must decide what to do with a question the detector cannot parse.

\paragraph{Procedure}
Each split draws three image-disjoint folds (fit, calibration, test); the score is a logistic regression with five
features (the VLM's probability, its confidence margin, the detector score, the detector score signed by the VLM's
answer, and the missing flag); thresholds and gate are fixed from the fit fold as in \S\ref{sec:fixing}; the
fixed-sequence test uses Clopper--Pearson bounds at $\delta=0.10$; and the deployed policy is evaluated on the test fold.
We use $400$ paired splits per cell, so every arm in a comparison sees identical folds.

\paragraph{Uncertainty}
Repeated random splits of one benchmark measure how sensitive the procedure is to the split. They do not
measure how the result would change on a fresh benchmark: the interval over splits shrinks like $1/\sqrt{\text{splits}}$ while
uncertainty about the dataset does not shrink at all. We therefore report gains as the mean over the $400$ splits
with a $95\%$ percentile interval from an \emph{outer image-level bootstrap}: $300$ bootstrap datasets
are drawn by resampling images with replacement, the full protocol is run on each with $20$ splits, and the
interval is taken over bootstrap datasets. A duplicated image keeps its identifier and is always placed in one fold. Because a
bootstrap dataset contains only about $63\%$ distinct images, the interval is somewhat wider than the one a
fresh sample of the same size would give. We do not report $p$-values over splits.

\paragraph{Validity audits}
Realised test-fold risk carries binomial noise and cannot audit a population guarantee: a policy whose true risk is
$0.08$ exceeds $0.10$ on a $150$-item test fold with substantial probability. We report two statistics per
cell. \emph{exc} is the fraction of non-aborted splits whose test-fold risk exceeds $\alpha$; under validity it
is bounded by $\delta$ plus finite-fold noise. \emph{cv} is the fraction of splits for which the one-sided
$90\%$ Clopper--Pearson \emph{lower} bound on the test-fold risk exceeds $\alpha$ (a confident violation); if the
population risk is at most $\alpha$ it is at most $0.1$ per split by the same argument as the theorem. Test folds
are used for nothing else: no score, threshold, gate or start is selected on them. An exact audit requires a known
population, which the simulation in \S\ref{sec:validity-audit} provides.
